\documentclass[12pt]{article}
\usepackage[utf8]{inputenc}
\usepackage[T1]{fontenc}
\usepackage[margin=1in]{geometry}
\usepackage{setspace}
\usepackage{xcolor}
\usepackage[colorlinks=true, urlcolor=blue!60!black, linkcolor=blue!60!black]{hyperref}
\setlength{\parindent}{0em}
\setlength{\parskip}{0.6em}
\renewcommand{\baselinestretch}{1.0}
\begin{document}
\begin{center}
{\LARGE \textit{Allen v.\ Wright}, 468 U.S.\ 737 (1984)\par}
\vspace{0.5em}
{\large Case Brief --- Two-Page Edition\par}
\end{center}

\section{Memory Jogger}
Black parents lacked standing to sue the IRS over lax enforcement against discriminatory private schools: stigma without personal denial of equal treatment is not cognizable, and desegregation harms were not ``fairly traceable'' to the IRS (\href{https://supreme.justia.com/cases/federal/us/468/737/}{Justia}).

\section{Detailed Case Facts}
The IRS denied tax exemptions to racially discriminatory private schools but required only that schools certify a nondiscriminatory policy. Parents of Black children in desegregating districts in seven States brought a nationwide class action claiming the rules let ``segregation academies'' keep exemptions. They never alleged their children had applied to any private school; they claimed stigma from federal aid and impairment of their children's chance at a desegregated public education.

\section{Procedural History}
The district court dismissed for lack of standing; a divided D.C.\ Circuit reversed and enjoined the IRS.

\section{Judicial Votes}
Reversed 5--3. O'Connor wrote for Burger, White, Powell, and Rehnquist; Brennan dissented, as did Stevens (joined by Blackmun). Marshall took no part.

\section{Holding}
Plaintiffs lack Article~III standing. Stigmatic injury counts only for those \textit{personally denied equal treatment}, and the desegregation injury, which depends on independent choices of private schools and white parents, is not fairly traceable to the IRS.

\section{Analysis of Opinions}
\textbf{Majority:} Standing rests on separation of powers. The desegregation injury was serious, but under \textit{Simon} the chain from tax exemptions to enrollment was speculative, and standing here would make courts ``continuing monitors'' of the Executive. \textbf{Brennan:} the Court used ``standing to slam the courthouse door'' and ignored the complaint's named schools. \textbf{Stevens:} causation followed from ``elementary economics,'' and the Court confused standing with the merits.

\section{Examples: Future Applications}
\textbf{Same-Side:} States suing to force more immigration arrests lack standing (\textit{United States v.\ Texas}, 2023). \textbf{Opposite-Side:} a Black student personally rejected by a tax-exempt academy because of her race has standing. \textbf{Fence-Sitter:} fuel producers challenging an electric-vehicle mandate whose injury runs through automakers and consumers, held traceable in \textit{Diamond Alternative Energy} (2025).

\section{Critique}
Scholars argue the Court manipulated causation to avoid a sensitive merits dispute and imposed a summary-judgment burden at the pleading stage. Progressives note that regulated entities can challenge enforcement while beneficiaries are locked out when agencies underenforce.

\section{Key Quotations}
\begin{quote}
``More important, the law of Art.\ III standing is built on a single basic idea --- the idea of separation of powers.''
\end{quote}
\noindent --- Justice O'Connor, for the Court, 468 U.S.\ at 752.

\end{document}
